\begin{proposition}[The logarithmic direct images by vertical twisting and duality]\hypertarget{tag-01RS}{}\label{inv:A1.2}\label{wb:prop:log-relative-direct-images}
\StatementRecord{inv:A1.2}{proved}{01RS}{\SeriesRef{inv:eq:A-H-images}, \SeriesRef{inv:eq:A-E-line}, \SeriesRef{inv:eq:A-N-images}, \SeriesRef{inv:A1.0}}

One has, as holomorphic vector bundles on $B$,
\[
\begin{array}{c|ccc}
q&0&1&2\\\hline
R^qf_*\mathcal E&\OO\oplus\OO(-1)&
 \OO^2\oplus\OO(-1)^2&\OO\oplus\OO(-1).
\end{array}
\]
They commute with base change to every full scheme fibre.
\end{proposition}


\begin{proof}
\begin{enumerate}[label=\arabic*.,leftmargin=*]
\item Put $H=2S_1+3S_2$.  Sections of the line bundle $\OO_X(H)$
on the cover at $p_j$ allow poles of order at most $m_j-1$
in $s$.  A fixed character has invariant exponent zero and gains no
pole, whereas the variable character $\zeta^{-a}$ has least allowed
invariant exponent $a-m$.  Thus the inclusion
$R^1f_*\OO\to R^1f_*\OO(H)$ is the identity on $\OO\gamma$ and
changes $s^ab_1$ to $s^{a-m}b_1=t_j^{-1}s^ab_1$ on its quotient.
It makes exactly one upper modification at each finite point.
The same computation applies to $R^2$.  No modification is made at
the cusp or a regular fibre.  The result is
\begin{equation}\hypertarget{tag-01T0}{}\label{inv:eq:A-H-images}
f_*\OO(H)=\OO,\qquad
R^1f_*\OO(H)=\OO\oplus\OO(1),\qquad
R^2f_*\OO(H)=\OO(1).
\end{equation}
Here the middle extension splits because $H^1(\OO(-1))=0$.
The local invariant modules are free.  Tensoring them with
$\C\{t_j\}/(t_j)$ gives invariants of
$s^{-(m_j-1)}\C\{s\}/(s^m)$, with the indicated surviving character
basis vectors.  The Reynolds idempotent is linear over the base,
so this calculation also proves the comparison with the nonreduced
fibre.  At the cusp the line bundle is trivial.  Cohomology and base
change, applied successively from the top degree, proves the claimed
local freeness and comparison for \eqref{inv:eq:A-H-images}.

\item The relative duality map we use is the cup product followed
by the analytic trace, giving a morphism of bundles
\[
R^qf_*(\omega_f\otimes\OO(-H))\longrightarrow
\bigl(R^{2-q}f_*\OO(H)\bigr)^\vee.
\]
Every fibre is a compact Gorenstein surface: it is the Cartier fibre
in a smooth threefold.  Serre duality on that scheme makes the map
an isomorphism on its residue field, also for $m_jS_j$.  The constant
dimensions just computed imply local freeness on both sides, so
Nakayama's lemma makes the bundle map an isomorphism.  Therefore
\begin{equation}\hypertarget{tag-01T1}{}\label{inv:eq:A-N-images}
(f_*N,R^1f_*N,R^2f_*N)
=(\OO(-1),\OO\oplus\OO(-1),\OO).
\end{equation}

\item The connecting maps of \eqref{inv:eq:A-E-line} are zero on a regular
fibre: its relative cotangent is a trivial bundle and evaluation on
a constant nonzero vector splits.  Source and target direct images
in \eqref{inv:eq:A-N-images} and \cref{inv:A1.0} are vector bundles,
so these holomorphic connecting maps, zero on a dense open set, are
zero everywhere.  We obtain extensions in each degree of the listed
$R^qN$ by $R^q\OO$.  Every resulting pair of line-bundle degrees is
zero or minus one; their extension groups are
$H^1(\OO(-1)),H^1(\OO),H^1(\OO(1))$, all zero.  These extensions
give exactly the three asserted bundles.  Applying the already proved
fibre comparisons to the short exact sequences proves base change
for $R^q\mathcal E$ as well.  
\end{enumerate}
\end{proof}